\scanpage{64}
% archive sheet number: 63
% @ L1–2
sur les circuits fondamentaux sur $\mathscr X$, qui sont de deux espèces

% @ L3–7, a)
a) \inserted{demi} circuit autour d'une \deleted{\unsure{droite|arête}}\inserted{composante} exceptionnelle de $K$ (en la traversant \deleted{une} \unsure{fois|fois ?}), soit $l_{\underline D_i}$ le circuit correspondant : la \unsure{psdroite|ps.\ droite} $D_i\in\Sigma$. On trouve
\[
u(l_{\vec D_i})=\text{translation par }n-1
\]
% @ L8–11
(dans le sens de l'orientation de la position rel.$^{\unsure{t|\acute e}}$ \unsure{à} $D$ voisine de $D_i$), valeurs des caractères fondamentaux
\[
\chi_X(l_{\underline D_i})=+1,\quad\chi_s(l_{\underline D_i})\underset{*}{=}\deleted{\unsure{+}}(-1)^{n-1},\quad\deleted{\chi_c(l_{\underline D_i})=1}
\]
% @ L12–17, b)
b) Demi-circuit autour d'un pt \deleted{de} $X$ (qui peut être soit dans une face, soit dans un côté, soit \inserted{être} un sommet). La \deleted{point} \unsure{feuillet|feuillets|facette|effet} de ce point dans $\mathrm{Pol}(D)$ reste \unsure{fixe|figée} pour $u(l)$, qui est de signe $-1$, donc on trouve toutes les symétries, donc
\[
u_D:\pi_1(\mathscr X,F_D)\longrightarrow\operatorname{Aut}(F_D)
\]
% @ L19
est \uline{surjectif} \quad ouf !

% @ R1–9, asterisk
$*$ Il semble que le caractère $\chi_c$ soit \uline{trivial} ! En effet, si pour une p.v. $D$ donnée, on désigne par $\mathrm{sing}_D$ (``singularité de $D$'') le nb de sommets de $\Sigma$ qu'elle contient, je dis que pour un chemin $c$ qui va de $D$ en $D'$, on a
\[
\mathrm{poids}\ c\equiv\mathrm{sing}_{D'}-\mathrm{sing}_D\qquad(2)
\]
% @ R10–15
(donc poids $c$ \uline{pair} si $D=D'$!) On est ramené au cas où $c$ est de longueur $1$, auquel cas cette formule est triviale. (en se rappelant que pour le franchissement \deleted{d'une d.} en ``balayage'' au dessus d'une \deleted{droite} $D_i\in\Sigma$, le poids est 2.

% @ R16–19
Donc finalement, \uline{on trouve}
\[
\boxed{\chi_{\mathscr X}=\chi_X\ -}
\]
il n'y a pas d'autre caractère à regarder que ceux provenant de $\mathbb D_2$ !
